% -------------------------
% Definition von Bedingungsflags
% -------------------------

\newif\ifcompany    % Definiert ein neues Bedingungsflag namens 'ifcompany'.
                    % Dieses Flag kann in einem LaTeX-Dokument verwendet werden, 
                    % um bestimmte Abschnitte des Textes basierend darauf ein- oder auszublenden.
                    % Standardmäßig ist das Flag 'false'.

\newif\ifreviewer   % Definiert ein weiteres Bedingungsflag namens 'ifreviewer'.
                    % Ähnlich wie bei 'ifcompany', kann dieses Flag verwendet werden,
                    % um Text zu steuern, der nur angezeigt wird, wenn 'ifreviewer' wahr ist.
                    % Standardmäßig ist dieses Flag ebenfalls 'false'.

\newif\ifsecondreviewer % Definiert ein drittes Bedingungsflag namens 'ifsecondreviewer'.
                    % Dieses Flag wird verwendet, um zu bestimmen, ob es einen zweiten Gutachter gibt.
                    % Standardmäßig ist das Flag 'false', was bedeutet, dass es keinen zweiten Gutachter gibt.

\newif\ifsupervisor % Definiert ein viertes Bedingungsflag namens 'ifsupervisor'.
                    % Dieses Flag ermöglicht es, Teile des Dokuments nur dann anzuzeigen,
                    % wenn 'ifsupervisor' auf 'true' gesetzt ist.
                    % Auch hier ist das Flag standardmäßig 'false'.

\newif\ifsecondauthor   % Erstellt ein Flag namens 'secondauthor'
                        % Das Flag wird verwendet, um zu bestimmen, ob es einen zweiten Autor gibt.
                        % Standardmäßig ist das Flag 'false', was bedeutet, dass es keinen zweiten Autor gibt.